% ID: SEL-D0053-OA-E02
% ESTADO: validada. S054.
\textbf{a)} $M(1)=28/14=\boxed{2}$. La ecuación $M(t)=5$ equivale a $11t+17=10t+60$, luego se alcanzan los cinco montajes diarios el día $\boxed{43}$.

\textbf{b)} Reescribimos $M(t)=11/2-49/(2t+12)$. Al crecer $t$ indefinidamente, el segundo término tiende a cero y $M(t)$ se aproxima a $5{,}5$ desde abajo. Ese valor es una asíntota horizontal del modelo.

\textbf{c)} $M'(t)=98/(2t+12)^2>0$ en todo el dominio. Por tanto, aumenta en el modelo continuo y también al comparar días enteros consecutivos.

\textbf{d)} La gráfica comienza en $(1,2)$ y pasa por $(43,5)$. Es creciente y se aplana al acercarse a la asíntota.
\begin{center}
\setlength{\unitlength}{0.9mm}
\begin{picture}(123,56)(-9,-7)
\qbezier(0,0)(48,0)(96,0)
\qbezier(0,0)(0,22)(0,44)
\put(97,0){$t$}
\put(1,45){$M(t)$}
{\color{gray}\qbezier(0,38.5)(48,38.5)(96,38.5)}
\qbezier(1.8,14.0000)(2.7,15.7354)(3.6,17.0625)
\qbezier(3.6,17.0625)(5.4,19.6826)(7.2,21.3500)
\qbezier(7.2,21.3500)(9.9,23.8451)(12.6,25.3077)
\qbezier(12.6,25.3077)(16.2,27.2736)(19.8,28.4118)
\qbezier(19.8,28.4118)(24.3,29.8521)(28.8,30.7045)
\qbezier(28.8,30.7045)(35.1,31.9036)(41.4,32.5862)
\qbezier(41.4,32.5862)(49.5,33.4747)(57.6,33.9868)
\qbezier(57.6,33.9868)(67.5,34.6215)(77.4,35.0000)
\qbezier(77.4,35.0000)(83.7,35.2479)(90,35.4375)
\put(99,37){\scriptsize $y=5{,}5$}
\put(4,10){\scriptsize $(1,2)$}
\put(74,29){\scriptsize $(43,5)$}
\end{picture}
\end{center}
